\section{Desarrollo}

El procedimiento se organizó en cinco fases. Cada una incluye los comandos ejecutados y
la evidencia gráfica del resultado obtenido.

% =====================================================================
\subsection{Fase 1: verificación del estado inicial}
% =====================================================================

Antes de instalar cualquier cosa conviene averiguar qué hay. Saltarse este paso es la
causa más común de terminar con dos o tres versiones del JDK conviviendo sin que el
usuario sepa cuál está usando realmente.

\begin{lstlisting}[language=bash,caption={Diagnóstico del estado inicial},morekeywords={java,javac,echo,which,readlink,dpkg,grep}]
# Version del entorno de ejecucion
java -version

# Version del compilador: si responde "command not found",
# lo que hay instalado es el JRE y no el JDK
javac -version

# Contenido de la variable de entorno: si sale vacio, no esta configurada
echo "JAVA_HOME = [$JAVA_HOME]"

# Ubicacion real del binario, siguiendo los enlaces simbolicos
which java
readlink -f $(which java)
\end{lstlisting}

En el equipo utilizado, \texttt{java -version} respondió con OpenJDK~11, mientras que
\texttt{javac -version} devolvió \texttt{command not found} y \texttt{JAVA\_HOME}
apareció vacía. Es decir: había un entorno de ejecución, pero no un kit de desarrollo.

\evidencia{01-verificacion-inicial.png}{Estado inicial del sistema: existe la JVM pero no el compilador, y \texttt{JAVA\_HOME} está vacía}{verificacion-inicial}

También conviene inventariar los paquetes relacionados que el gestor tiene registrados,
porque son los que habrá que retirar:

\begin{lstlisting}[language=bash,caption={Inventario de paquetes Java instalados},morekeywords={dpkg,grep,awk,ls}]
# Paquetes de Java registrados en el sistema
dpkg -l | grep -iE "jdk|jre"

# Versiones presentes en el directorio estandar de JVM
ls -la /usr/lib/jvm/
\end{lstlisting}

\evidencia{02-paquetes-instalados.png}{Paquetes relacionados con Java presentes antes de la instalación}{paquetes-instalados}

% =====================================================================
\subsection{Fase 2: desinstalación de versiones previas}
% =====================================================================

Esta fase es opcional en un equipo limpio, pero es la que evita el escenario más
confuso de todos: tener varias versiones instaladas y no saber cuál responde.

\begin{lstlisting}[language=bash,caption={Retiro de la versión anterior},morekeywords={sudo,apt,remove,purge,autoremove,java,javac}]
# --purge elimina tambien los archivos de configuracion, no solo el paquete
sudo apt remove --purge "openjdk-11-*" -y

# Retira las dependencias que quedaron huerfanas
sudo apt autoremove -y

# Comprobacion: ambos comandos deben responder "command not found"
java -version
javac -version
\end{lstlisting}

\evidencia{03-desinstalacion.png}{Salida del proceso de desinstalación de OpenJDK 11}{desinstalacion}

\evidencia{04-verificacion-limpieza.png}{Verificación de que el sistema quedó sin ninguna versión de Java}{verificacion-limpieza}

\textbf{Incidencia registrada durante esta práctica.} Al ejecutar el primer comando,
el sistema respondió:

\begin{lstlisting}[language=bash,caption={Error por una operación de dpkg interrumpida},morekeywords={sudo,apt,remove,purge,dpkg,configure}]
$ sudo apt remove --purge "openjdk-11-*" -y
E: se interrumpio la ejecucion de dpkg, debe ejecutar manualmente
   «sudo dpkg --configure -a» para corregir el problema
\end{lstlisting}

El mensaje no tiene relación con Java: indica que una instalación anterior de
\textbf{cualquier} paquete quedó a medio configurar, y APT se niega a realizar
operaciones nuevas mientras la base de datos de paquetes esté en ese estado
inconsistente. La reparación es la que el propio mensaje sugiere:

\begin{lstlisting}[language=bash,caption={Reparación del estado de dpkg},morekeywords={sudo,dpkg,configure,apt,fix-broken,install,update,audit}]
# Termina la configuracion pendiente
sudo dpkg --configure -a

# Repara dependencias rotas, si las hubiera
sudo apt --fix-broken install

# Comprueba que APT vuelve a operar con normalidad
sudo apt update
\end{lstlisting}

Si el problema persiste, \texttt{sudo dpkg --audit} enumera los paquetes que quedaron
en estado inconsistente, lo que permite identificar cuál es el que falla.

\evidencia{03a-dpkg-interrumpido.png}{Error por una operación de \texttt{dpkg} interrumpida y su reparación}{dpkg-interrumpido}

\textbf{Advertencia importante.} En Ubuntu, varias aplicaciones de escritorio dependen
de Java. Si \texttt{apt autoremove} propone eliminar programas que se están usando,
conviene revisar la lista antes de aceptar. Cuando existan dudas, la alternativa
prudente es \textbf{no desinstalar}: instalar la versión nueva junto a la anterior y
seleccionar cuál usar con \texttt{update-alternatives}, como se explica en la fase~4.

% =====================================================================
\subsection{Fase 3: descarga e instalación del JDK}
% =====================================================================

\subsubsection{Actualización de los índices de paquetes}

APT trabaja con una copia local del catálogo de los repositorios. Si esa copia está
desactualizada, el sistema puede no encontrar el paquete o intentar descargar una
versión que ya no existe en el servidor —un error frecuente que se manifiesta como
\texttt{404 Not Found} a mitad de la instalación—.

\begin{lstlisting}[language=bash,caption={Actualización del catálogo de paquetes},morekeywords={sudo,apt,update}]
sudo apt update
\end{lstlisting}

\evidencia{05-actualizar-indices.png}{Actualización de los índices de los repositorios}{actualizar-indices}

\subsubsection{Consulta de la versión disponible}

Antes de instalar conviene confirmar qué versión exacta ofrece el repositorio y de
cuál procede:

\begin{lstlisting}[language=bash,caption={Consulta de la versión candidata},morekeywords={apt-cache,policy,apt,search}]
# Version candidata y repositorio de origen
apt-cache policy openjdk-21-jdk

# Todas las versiones del JDK disponibles en los repositorios
apt search --names-only "^openjdk-.*-jdk$"
\end{lstlisting}

En Ubuntu 22.04 el paquete proviene del repositorio \textit{jammy-updates}, en la
sección \textit{universe}, con la versión \texttt{21.0.11+10-1\textasciitilde 22.04.2}.

\evidencia{06-version-disponible.png}{Versión candidata de \texttt{openjdk-21-jdk} y repositorio del que proviene}{version-disponible}

\subsubsection{Instalación}

\begin{lstlisting}[language=bash,caption={Instalación del JDK},morekeywords={sudo,apt,install}]
sudo apt install openjdk-21-jdk -y
\end{lstlisting}

Nótese que el paquete es \texttt{openjdk-21-\textbf{jdk}} y no
\texttt{openjdk-21-\textbf{jre}}. Instalar el segundo por descuido deja el sistema en el
mismo estado del que se partió: con máquina virtual pero sin compilador.

\evidencia{07-instalacion.png}{Proceso de instalación de OpenJDK 21}{instalacion}

\subsubsection{Verificación de la instalación}

\begin{lstlisting}[language=bash,caption={Comprobación de que el JDK quedó instalado},morekeywords={java,javac,version}]
java -version      # entorno de ejecucion
javac -version     # compilador: ahora si debe responder
\end{lstlisting}

Ambos comandos deben reportar la versión 21. Que \texttt{javac} responda es la señal de
que se instaló un JDK y no solo un JRE.

\evidencia{08-verificacion-version.png}{Versiones de \texttt{java} y \texttt{javac} tras la instalación}{verificacion-version}

% =====================================================================
\subsection{Fase 4: configuración de las variables de entorno}
% =====================================================================

Esta es la parte central de la práctica y la que más errores concentra.

\subsubsection{Identificación de la ruta de instalación}

\texttt{JAVA\_HOME} debe apuntar al \textbf{directorio raíz} del JDK: aquel que contiene
las carpetas \texttt{bin}, \texttt{lib} y \texttt{conf}. No al ejecutable, ni al
subdirectorio \texttt{bin}.

\begin{lstlisting}[language=bash,caption={Localización del directorio raíz del JDK},morekeywords={readlink,which,dirname,ls,update-alternatives}]
# Ruta real del compilador, resolviendo los enlaces simbolicos.
# Devuelve algo como /usr/lib/jvm/java-21-openjdk-amd64/bin/javac
readlink -f $(which javac)

# Subir dos niveles desde el binario da el directorio raiz del JDK
dirname $(dirname $(readlink -f $(which javac)))

# Contenido del directorio: deben aparecer bin, lib y conf
ls /usr/lib/jvm/
ls /usr/lib/jvm/java-21-openjdk-amd64/
\end{lstlisting}

En Ubuntu la ruta resultante es \texttt{/usr/lib/jvm/java-21-openjdk-amd64}.

\evidencia{09-ruta-jdk.png}{Identificación del directorio raíz de la instalación del JDK}{ruta-jdk}

\subsubsection{Definición de las variables}

Hay que decidir el alcance de la configuración, y la elección tiene consecuencias:

\rowcolors{2}{rowblue}{white}
\footnotesize
\renewcommand{\arraystretch}{1.15}
\setlength{\tabcolsep}{4pt}

\begin{longtable}{|>{\raggedright\arraybackslash}p{3.2cm}|>{\raggedright\arraybackslash}p{2.6cm}|>{\raggedright\arraybackslash}p{8.4cm}|}
	\caption{Archivos donde puede definirse \texttt{JAVA\_HOME} y su alcance}
	\label{tab:alcance} \\
	\hline
	\rowcolor{headerblue}
	\headercell{Archivo} & \headercell{Alcance} & \headercell{Cuándo conviene} \\
	\hline
	\endfirsthead
	\hline
	\rowcolor{headerblue}
	\headercell{Archivo} & \headercell{Alcance} & \headercell{Cuándo conviene} \\
	\hline
	\endhead

	\texttt{\textasciitilde/.bashrc} & Un usuario, solo en terminales interactivas &
	Es el caso más común en un equipo personal. No lo leen los servicios del sistema ni,
	en algunas configuraciones, las aplicaciones lanzadas desde el menú gráfico. \\ \hline
	\texttt{\textasciitilde/.profile} & Un usuario, también en la sesión gráfica &
	Cuando se necesita que los programas abiertos desde el escritorio —un IDE, por
	ejemplo— también vean la variable. \\ \hline
	\texttt{/etc/environment} & Todo el sistema & Equipos con varios usuarios. No admite
	sustituciones ni la sintaxis \texttt{export}: solo pares \texttt{CLAVE=valor}. \\ \hline
	\texttt{/etc/profile.d/*.sh} & Todo el sistema & La opción más limpia para el
	alcance global: es un archivo aparte y no se pierde al actualizar el sistema. \\ \hline
\end{longtable}

\normalsize
\setlength{\tabcolsep}{6pt}

En este reporte se configuró a nivel de usuario, editando \texttt{\textasciitilde/.bashrc}:

\begin{lstlisting}[language=bash,caption={Definición de \texttt{JAVA\_HOME} y ampliación del \texttt{PATH}},morekeywords={nano,export,cat,tee,source}]
nano ~/.bashrc
\end{lstlisting}

Y agregando al final del archivo:

\begin{lstlisting}[language=bash,caption={Líneas añadidas a \textasciitilde/.bashrc},morekeywords={export}]
# --- Configuracion de Java ---
export JAVA_HOME=/usr/lib/jvm/java-21-openjdk-amd64
export PATH=$JAVA_HOME/bin:$PATH
\end{lstlisting}

Dos observaciones sobre estas dos líneas:

\begin{itemize}
	\item \texttt{PATH} se \textbf{amplía}, no se reemplaza. La parte \texttt{:\$PATH}
	al final conserva todo lo que ya había. Omitirla deja la sesión sin acceso a
	\texttt{ls}, \texttt{cd} ni ningún otro comando, y es uno de los errores más
	desconcertantes que pueden cometerse en este paso.
	\item El directorio del JDK se coloca \textbf{al principio} del \texttt{PATH}, de
	modo que si hubiera otra versión de Java en el sistema, esta tenga precedencia.
\end{itemize}

\evidencia{10-editar-bashrc.png}{Líneas agregadas al final de \texttt{\textasciitilde/.bashrc}}{editar-bashrc}

\subsubsection{Aplicación y verificación}

Un archivo de configuración editado no surte efecto en la terminal que ya estaba
abierta: hay que releerlo o abrir una terminal nueva.

\begin{lstlisting}[language=bash,caption={Aplicación y comprobación de las variables},morekeywords={source,echo,ls}]
# Relee el archivo en la terminal actual
source ~/.bashrc

# La variable debe mostrar la ruta del JDK
echo $JAVA_HOME

# El directorio del JDK debe aparecer al inicio del PATH
echo $PATH

# Prueba definitiva: invocar el binario a traves de la variable
$JAVA_HOME/bin/java -version
$JAVA_HOME/bin/javac -version
\end{lstlisting}

La última comprobación es la más significativa del reporte: si
\texttt{\$JAVA\_HOME/bin/java} responde con la versión, entonces la variable apunta a un
directorio válido. Si responde \texttt{No such file or directory}, la ruta está mal
escrita o apunta a un nivel equivocado del árbol de directorios.

\evidencia{11-java-home-configurado.png}{Variables \texttt{JAVA\_HOME} y \texttt{PATH} configuradas y verificadas}{java-home-configurado}

\subsubsection{Gestión de varias versiones instaladas}

Cuando conviven varios JDK, Debian y Ubuntu resuelven cuál responde mediante el sistema
de alternativas:

\begin{lstlisting}[language=bash,caption={Selección de la versión activa},morekeywords={sudo,update-alternatives,config,list,java,javac}]
# Versiones registradas en el sistema
update-alternatives --list java

# Menu interactivo para elegir cual sera la predeterminada
sudo update-alternatives --config java
sudo update-alternatives --config javac
\end{lstlisting}

Conviene entender que \texttt{update-alternatives} y \texttt{JAVA\_HOME} son mecanismos
\textbf{independientes}: el primero decide qué binario responde al escribir
\texttt{java} en la terminal; el segundo, qué JDK usan Maven, Gradle o el IDE. Pueden
apuntar a versiones distintas sin que el sistema se queje, y esa discrepancia silenciosa
explica buena parte de los problemas de compilación difíciles de diagnosticar.

\evidencia{12-alternatives.png}{Versiones registradas en el sistema de alternativas}{alternatives}

% =====================================================================
\subsection{Fase 5: verificación funcional}
% =====================================================================

Que los comandos reporten una versión no demuestra que el kit funcione. La prueba real
es compilar y ejecutar un programa.

\begin{lstlisting}[language=bash,caption={Creación del programa de prueba},morekeywords={mkdir,cd,cat,nano}]
mkdir -p ~/pruebaJava && cd ~/pruebaJava
nano HolaMundo.java
\end{lstlisting}

\begin{lstlisting}[caption={Programa de prueba (\texttt{HolaMundo.java})}]
public class HolaMundo {
    public static void main(String[] args) {
        System.out.println("El JDK quedo instalado correctamente.");
        System.out.println("Version de Java : " + System.getProperty("java.version"));
        System.out.println("Proveedor       : " + System.getProperty("java.vendor"));
        System.out.println("Directorio home : " + System.getProperty("java.home"));
    }
}
\end{lstlisting}

El programa no se limita a imprimir un saludo: consulta las propiedades del sistema para
que la salida muestre qué JDK está ejecutándolo realmente. El valor de
\texttt{java.home} debe coincidir con el de \texttt{JAVA\_HOME}; si no coincide, algo de
la configuración quedó inconsistente.

\evidencia{13-holamundo-codigo.png}{Código fuente del programa de prueba}{holamundo-codigo}

\begin{lstlisting}[language=bash,caption={Compilación y ejecución},morekeywords={javac,java,ls}]
# Compilar: genera HolaMundo.class
javac HolaMundo.java

# Verificar que se creo el bytecode
ls -l

# Ejecutar (sin la extension .class)
java HolaMundo
\end{lstlisting}

\evidencia{14-holamundo-ejecucion.png}{Compilación y ejecución del programa de prueba}{holamundo-ejecucion}

Con esta última evidencia queda demostrado el ciclo completo: el código fuente se
compiló a \textit{bytecode} con \texttt{javac} y la JVM lo ejecutó, usando el JDK
señalado por \texttt{JAVA\_HOME}.
